% Printed 391, beginning at the dated heading; preceding header is in the prior component.
\begin{center}
\textsc{Gauss} an \textsc{Gerling}.\quad Göttingen, 12.~September 1838.
\end{center}
\par\noindent\hspace*{\parindent}\mbox{.\ .\ .\ .\ .\ .}\quad Was den Längenunterschied mit Paris betrifft, so habe ich selbst die Resultate aus Sternbedeckungen etc., die \textsc{Triesnecker}, \textsc{Wurm} u.~a. gezogen haben, niemals zusammengestellt: Sie selbst sind also ebenso gut wie ich selbst im Stande, \textls[120]{dieses} Resultat zu ermitteln. Dagegen habe ich aus den Längenunterschieden zwischen Jever und Göttingen, imgleichen Bentheim und Göttingen, wie Sie aus meinen, resp. meines Sohnes, trigonometrischen Messungen folgten, einerseits, und Jever und Bentheim mit Paris andererseits wie sie \textsc{Krayenhoff} ansetzt, denjenigen Längenunterschied zwischen Göttingen und Paris abgeleitet, welchen Sie in \textsc{Hardings} Ephemeriden angesetzt finden. Eine neue Rechnung habe ich deshalb nicht gemacht; bei einer completen Ausgleichung aller Dreiecke, die als Kranz das Oldenburgsche umgeben, leidet Jever eine, doch nur sehr geringe Abänderung; auch könnte Jever ganz weggelassen und dagegen Emden und Onstwedde gebraucht werden, bis wohin \textsc{Hartmanns} Messungen sich erstreckt haben. Indessen würde jedenfalls die Abänderung des Endresultats nur ganz unbedeutend sein, und ich habe es bisher um so weniger der Mühe werth gehalten, deshalb eine neue Rechnung zu machen, weil \textsc{Krayenhoffs} Messungen doch, wie Sie aus meinem Supplem. theor. comb. abnehmen können, weit schlechter sind als ihr Ruf.
\par\noindent\hspace*{\parindent}Eine dritte Bestimmung könnten Sie aus Altona entnehmen, dessen Differenz gegen Göttingen trigonometrisch bestimmt ist (siehe Breitenunterschied [Art.~19]), sowie die Differenz von Greenwich chronometrisch. Aber von mir selbst ist das Betreffende nicht zusammengestellt, was ich daher Ihnen überlassen muss.\quad\mbox{.\ .\ .\ .\ .\ .}

% Printed 392.
\par\clearpage
\noindent\makebox[0pt][l]{\fontsize{13}{16}\selectfont 392}\hfill{\small\textbf{BRIEFWECHSEL MIT GERLING.}}\hfill\null\par\medskip
\begin{center}
\textsc{Gauss} an \textsc{Gerling}.\quad Göttingen, 14.~November 1838.
\end{center}
\par\noindent\hspace*{\parindent}\mbox{.\ .\ .\ .\ .\ .}\quad Dass ich auf Ihren vorletzten Brief noch nicht weiter geantwortet habe, müssen Sie damit gütigst entschuldigen, dass ich eigentlich zu dem in meinem letzten Briefe [vom 12.~September] mitgetheilten kaum noch etwas hinzu zu setzen wusste.
\par\noindent\hspace*{\parindent}Thatsächliches kann ich wirklich nichts beifügen, als die Hinweisung auf \textsc{Lindenau-Bohnenbergers} Zeitschrift, Band IV, S.~119, wo Sie die relative Lage der alten und neuen hiesigen Sternwarte bereits angesetzt finden. Allenfalls mit folgenden beiden Anmerkungen.
\par\noindent\hspace*{\parindent}1) Der Längenunterschied ist in \textls[120]{Zeit} angesetzt.
\par\noindent\hspace*{\parindent}2) Die Position der neuen Sternwarte a.~a.~O. bezieht sich auf das Centrum der Rotunde, dessen Lage gegen die Mitte der Axe des \textsc{Reichenbach}schen Meridiankreises ich Ihnen, wenn ich nicht irre, bereits früher mitgetheilt habe. Wollen Sie aber die relative Lage des letztern Punktes gegen das Centrum der alten Sternwarte selbst scharf berechnen, so schreibe ich dazu aus meinem Verzeichnisse der Coordinaten folgendes ab:
\begin{center}
\setlength{\tabcolsep}{4pt}
\begin{tabular}{@{}lrr@{}}
Centrum der Kuppel & \(-3{,}104\) & \(-7{,}324\)\\
Mitte der Axe des \textsc{Reichenbach}schen Meridiankreises & \(0\) & \(0\)\\
Mitte der alten Sternwarte & \(-193{,}54\) & \(+541{,}9\).
\end{tabular}
\end{center}
\par\noindent\hspace*{\parindent}Die Einheit ist das Meter; bei der ersten Zahl bedeutet \(+\) südlich, bei der zweiten \(+\) westlich.
\par\noindent\hspace*{\parindent}Verlangen Sie nun \textls[120]{meine} Meinung über die \textls[120]{Reductionszahl}, die Sie bei dem Ansetzen Ihrer Längen gegen Paris anzuwenden haben, so kann ich nur sagen, dass ich dies für etwas sehr gleichgültiges ansehe, wenn Sie mir zugleich mittheilen, \textls[120]{welche} Zahl Sie zum Grunde gelegt haben. In der That, Sie mögen wählen was Sie wollen, so bleibt die Latitüde in dieser Beziehung oder die Grösse des Spielraums
\par\noindent\hspace*{\parindent}\hangindent=2.8em\hangafter=1 A.\quad \textls[120]{höchst} unbedeutend für jeden praktischen Zweck, z.~B. Kartenzeichnungen.
\par\noindent\hspace*{\parindent}\hangindent=2.8em\hangafter=1 B.\quad zwanzig- oder fünfzigmal grösser als alle \textls[120]{relativen} Differenzen zwischen allen Ihren Ansätzen.

% Printed 393.
\par\clearpage
\noindent\null\hfill{\small\textbf{ZUR HANNOVERSCHEN TRIANGULATION.}}\hfill\makebox[0pt][r]{\fontsize{13}{16}\selectfont 393}\par\medskip
\par\noindent\hspace*{\parindent}Indem nun die letztern eigentlich das \textls[120]{Wesen Ihrer Arbeit} ausmachen, so scheint mir dadurch mein oben ausgesprochenes Urtheil hinlänglich gerechtfertigt.
\par\noindent\hspace*{\parindent}In Ihrem vorletzten Briefe sprechen Sie von den Differenzen des Längenunterschiedes zwischen Paris und Göttingen auf verschiedenen Wegen, die bis auf \(\frac{3}{4}\) Secunden in Zeit gehen, als sehr beträchtlichen. Ich würde sie nicht so nennen.
\par\noindent\hspace*{\parindent}In der That scheinen Sie mir dabei ganz ausser Acht gelassen zu haben, was ich in meiner Schrift über den Breitenunterschied zwischen Göttingen und Altona S.~73 [*)] entwickelt habe. Wenn man zwischen zwei Punkten, die nicht gar zu weit aus einander liegen, auf geodätischem Wege vermittelst zweier verschiedener Touren Längenunterschiede fände, die um 11 Bogensecunden verschieden wären, so wäre dies ein ganz ungeheurer Fehler, der bewiese, dass ganz enorm schlechte Beobachtungen untergelaufen wären, oder schlecht gerechnet.
\par\noindent\hspace*{\parindent}Aber Paris und Göttingen liegen doch sehr weit aus einander. \textsc{Delambres} und \textsc{Krayenhoffs} Messungen sind, glaube ich, an Genauigkeit mit den unsrigen gar nicht zu vergleichen, auch nach ganz verschiedenen Methoden und Elementen berechnet. Ich würde mich daher nicht wundern, wenn der so gefundene Längenunterschied um eine Anzahl Bogensecunden von einem andern auch geodätisch gefundenen, aber z.~B. über Strassburg, Mannheim, etc. geleiteten, abwiche.
\par\noindent\hspace*{\parindent}Allein \textls[120]{davon} ist ja hier gar nicht die Rede, sondern von dem Unterschiede der astronomisch bestimmten Längen, die ja, genau besehen, etwas \textls[120]{ganz anderes} bedeuten, als die geodätischen. Jene beziehen sich auf das Fortschreiten der Lage der Verticallinien und der durch sie und parallel mit der Erdaxe gelegten Plana, die andern auf die Distanzen auf der Oberfläche der Erde, die \textls[120]{gar kein Ellipsoid} ist, sondern wozu ein Ellipsoid nur wie eine Art Annäherung betrachtet werden kann.
\par\noindent\hspace*{\parindent}Nach so unzähligen vorliegenden Ungleichmässigkeiten des Geodätischen und Astronomischen in Beziehung auf die Breite muss man auf Differenzen bei der Länge von derselben Ordnung überall gefasst sein. Eine Differenz
\par\vfill
\noindent\rule{.22\linewidth}{.4pt}\par\smallskip
{\footnotesize\noindent\hspace*{\parindent}[*) S.~49/50 dieses Bandes.]\par}
\medskip\noindent{\small IX.}\hfill 50\hspace*{1.5em}\par

% Printed 394.
\clearpage
\noindent\makebox[0pt][l]{\fontsize{13}{16}\selectfont 394}\hfill{\footnotesize\textbf{BRIEFWECHSEL MIT GERLING,\quad ZUR HANNOVERSCHEN TRIANGULATION.}}\hfill\null\par\medskip
\noindent von \(11^{\prime\prime}\) selbst zwischen zwei Örtern, die einander \textls[120]{viel} näher lägen als Göttingen und Paris, ist im Grunde gar nichts. Es hätte schlechterdings nichts besonderes, wenn selbst eine doppelt so grosse sich zwischen Göttingen und Marburg fände.
\par\noindent\hspace*{\parindent}Übersehen Sie also nicht, dass die von Ihnen gefundenen Unterschiede sogar sehr gut selbst dann vollkommen genau richtig sein könnten, wenn \textls[120]{alle} Beobachtungen vollkommen und absolut fehlerfrei wären (was sie nicht sind). Es ist \textls[120]{gar kein} Grund zu erwarten, dass
\par\noindent\hspace*{\parindent}1) der durch geodätische Messungen, unter Voraussetzung bestimmter Ellipsoidsdimensionen, berechnete Unterschied zwischen Göttingen und Paris,
\par\noindent\hspace*{\parindent}2) der durch astronomische Beobachtungen gefundene,
\par\noindent\hspace*{\parindent}3) der gemischt gefundene, nemlich geodätisch zwischen Göttingen und Altona und astronomisch zwischen Altona-Greenwich-Paris,
\par\noindent unter sich übereinstimmen sollten, da wirklich diese drei Zahlen ganz verschiedene Dinge bedeuten.
\par\noindent\hspace*{\parindent}Wenn ich bei einer künftigen Bekanntmachung der Resultate meiner Messungen einen Theil davon in der Form von Breiten- und von Längendifferenzen geben sollte, so werde ich letztere ganz gewiss nicht in Beziehung auf Paris ansetzen, welches mich gar nichts angeht. Finden Sie aber aus was immer für Gründen es für gerathen, die Ihrigen in einer solchen Form mitzutheilen, so wiederhole ich, dass ich es für ganz gleichgültig halte, welche Reductionszahl Sie brauchen, wenn nur der Leser nicht in Ungewissheit bleibt, dass die dabei befindlichen Bruchtheile der Secunden an sich gar keinen Anspruch auf eine scharfe Bedeutung machen können, sondern nur als Mittel dienen, jede beliebige Differenz zwischen je zweien Ihrer Zahlen mit derjenigen Schärfe wieder zu erhalten, dass nichts von der Schärfe Ihrer Messungen an sich verloren geht. Ohne Ihnen vorzugreifen, würde ich unter der eben ausgesprochenen Voraussetzung diejenige Zahl zum Grunde legen, die taliter qualiter aus der vorhandenen geodätischen Verbindung zwischen Göttingen und Paris hervorgegangen ist, da eigentlich die etwas ganz anderes bedeutenden astronomischen Längen hier ganz ausser Frage sind.\quad\mbox{.\ .\ .\ .\ .\ .}
\begin{center}\rule{.12\linewidth}{.4pt}\end{center}

% Printed 395.
\par\clearpage
\noindent\null\hfill{\footnotesize\textbf{BEMERKUNGEN. BRIEFWECHSEL ZUR HANNOVERSCHEN TRIANGULATION.}}\hfill\makebox[0pt][r]{\fontsize{13}{16}\selectfont 395}\par\medskip
\begin{center}\textbf{BEMERKUNGEN.}\end{center}
\begingroup\small
\par\noindent\hspace*{\parindent}Der Abdruck der vorstehenden Briefe ist nach den Originalen erfolgt.
\par\noindent\hspace*{\parindent}Dem Brief an \textsc{Bessel} vom 26.~December 1821, S.~349/351, war eine Dreiecksskizze zugefügt, welche die 5 Dreiecke von der Seite Hohehagen-Inselsberg bis zur Seite Hils-Lichtenberg (vergl. die Figur auf S.~299) enthält, und ferner die Schnitte vom Hils aus nach dem Deister und dem Brelingerberge angibt. Der Brelingerberg, nördlich von Hannover, und der Wohlenberg, nordwestlich von Braunschweig, waren 1821 bei der Recognoscirung als Dreieckspunkte in Aussicht genommen. Der erstere wurde jedoch beim Beginn der Arbeiten im folgenden Jahre als untauglich erkannt, der zweite als überflüssig fortgelassen. Wahrscheinlich hat auch zu dem Brief an \textsc{Bessel} vom 15.~November 1822, S.~351/358, eine Karte gehört.
\par\noindent\hspace*{\parindent}Zu den Briefen an \textsc{Olbers} vom 4. und 8.~Juli 1824, S.~369/372, ist noch zu bemerken, dass der in ihnen erwähnte Punkt Wentel nicht in das System aufgenommen ist.
\par\noindent\hspace*{\parindent}Es war, wie es in dem \textsc{Gauss}schen Arbeitsbericht für 1824 an das hannoversche Ministerium heisst, »schneller als man bei den grossen Schwierigkeiten hatte hoffen können, ein sehr gutes Dreieckssystem bis Bremen gebildet.« »Allein desto trüber wurden nun die Aussichten für die weitere Fortsetzung der Dreiecke.\ \mbox{.\ .\ .\ .}\quad Nach dem Fehlschlagen aller Versuche, einen neuen Dreieckspunkt an Brüttendorf zu knüpfen, blieb noch ein Behelf übrig, nemlich den Steinberg mit in das Dreieckssystem aufzunehmen.« Aber auch an die Seite Bremen-Steinberg liess sich kein neues Dreieck schliessen. »Ohne einen eben jetzt eingetretenen glücklichen Umstand würde es um die Fortsetzung der Messungen sehr misslich gestanden haben, da alle Möglichkeiten jetzt erschöpft schienen.« Der erwähnte glückliche Umstand war die von \textsc{Gauss} auf dem Ansgariusthurm in Bremen gemachte Entdeckung, dass die Spitze des Thurmes von Zeven dort noch eben zu sehen war. Die Sichtbarkeit Bremens in Zeven war früher wegen Moorrauch nicht bemerkt worden; Zeven liess sich auch ausser mit Bremen mit den vorhergehenden Punkten Steinberg, Wilsede und Litberg verbinden. Von Bremen-Zeven konnte dann die Dreieckskette fortgesetzt werden. (Vergl. dazu den später folgenden Auszug aus dem \textsc{Gauss}schen Arbeitsbericht für 1825, sowie den gleichfalls noch folgenden Brief an \textsc{Olbers} aus dem Juli 1825.)
\par\noindent\hspace*{\parindent}Die im Briefe an \textsc{Gerling} vom 11.~August 1823, S.~385/386, erwähnten durch den Heliotrop übermittelten Zahlenzeichen haben (nach einigen kurzen Notizen in Beobachtungs- und Rechnungsheften zur Gradmessung) vorher verabredete, zum Theil von Station zu Station wechselnde Bedeutung gehabt (vergl. auch den Anfang des Briefes an \textsc{Gerling} vom 1.~September 1823, S.~387); z.~B. bedeutete nach einer Aufzeichnung 1: es geht so gut, 2: Licht soll geschwächt werden, 3: für heute soll aufgehört werden, 1\,.\,1: Abgang zum nächsten Punkt.
\par\noindent\hspace*{\parindent}Auf die Gradmessungsarbeiten in den Jahren 1821/1825 beziehen sich auch die Briefe an \textsc{Schumacher} vom 6. und 30.~Mai, 11.~Juli, 29.~September und 24.~October 1821; vom 10.~Mai, 10.~Juni, 6. und 30.~August, 6., 18., 24. und 29.~September, 8.~October und 10.~November 1822; vom 8. und 18.~Juni, 23.~Juli, Anfang August, 21.~August, 18.~September und 23.~October 1823; vom 24.~Juni, 1.~Juli, 27.~September, 17.~October, 28.~November 1824 und vom Januar 1825; vom 29.~April, 14., 20. und 26.~Juni, 11. und 15.~Juli und 14.~August 1825. (Vergl. den Briefwechsel \textsc{Gauss-Schumacher}, Erster Band, S.~229/245, S.~265/291, S.~311/325, S.~334/336, S.~397/413, S.~426/430. Zweiter Band, S.~1/3, S.~14/31.)
\par\noindent\hspace*{\parindent}Die mehrfach erwähnte Triangulation des Kurfürstenthums Hannover durch französische Ingenieurgeographen unter Leitung des Oberstlieutenants \textsc{Epailly} hat 1804 und 1805 während der französischen Occupation stattgefunden; sie sollte die Unterlagen für die kartographische Aufnahme geben. Das Tableau der \textsc{Epailly}schen Dreiecke, nebst einem Berichte \textsc{Epaillys} über seine Arbeiten, erhielt \textsc{Gauss} im Anfang des Jahres 1821 durch \textsc{Laplaces} Vermittelung vom französischen Kriegsministerium.
\par\noindent\hfill\textsc{Krüger}.\par
\endgroup
\begin{center}\rule{.12\linewidth}{.4pt}\end{center}
\vfill\noindent\hfill 50\textsuperscript{*}\hspace*{1.5em}\par

% Printed 396.
\clearpage
\noindent\makebox[0pt][l]{\fontsize{13}{16}\selectfont 396}\hfill{\small\textbf{ZUR HANNOVERSCHEN TRIANGULATION.}}\hfill\null\par\medskip
\begin{center}
{\large\textbf{VERÖFFENTLICHUNGEN}}\par\bigskip
{\large\textbf{ZUR HANNOVERSCHEN TRIANGULATION.}}\par\medskip
[II.]\par\bigskip
\rule{.74\linewidth}{.4pt}\\[-.5ex]
{\small Astronomische Nachrichten, Band I, Nr.~7, Februar 1822, S.~105---106.}\\[-.5ex]
\rule{.74\linewidth}{.4pt}
\end{center}
\par\noindent\hspace*{\parindent}\{Es ist bekannt, dass die hannoversche Regierung die von Sr.~Majestät dem Könige von Dänemark, diesem erhabenen Beförderer der Wissenschaften, begonnene Gradmessung, die sich in der Breite vom nördlichsten Punkte Jütlands bis zur südlichsten Grenze von Lauenburg erstreckt, fortsetzt, und die Ausführung dieser Arbeit dem Herrn Hofrath \textsc{Gauss} in Göttingen übertragen hat. Unsere gemeinschaftlichen Dreiecke sind Hamburg-Hohenhorn-Lüneburg und Hohenhorn-Lauenburg-Lüneburg. In diesen Dreiecken ist der Winkel in Hamburg von mir allein, die Winkel in Lüneburg von Herrn Hofrath \textsc{Gauss} und mir gemeinschaftlich beobachtet (der Unterschied unserer Messungen war unter \(0^{\prime\prime}{,}3\)); die Winkel in Hohenhorn und Lauenburg sind zu derselben Zeit (1818) von Herrn Capitain v.~\textsc{Caroc} gemessen. Im vorigen Jahre hat der Herr Hofrath seine Dreiecke an der südlichen Grenze begonnen, und der Auszug seines Briefes, den ich hier mittheile, enthält eine kurze Nachricht darüber.
\par\noindent\hfill\textit{S}[\textit{chumacher}].\}\par
\begin{center}Auszug aus einem Briefe des Herrn Hofrath und Ritter \textsc{Gauss}.\end{center}
\par\noindent\hspace*{\parindent}Den Zustand meiner Triangulation habe ich das Vergnügen auf beiliegenden Kärtchen Ihnen mitzutheilen. Die starken Linien sind die, wo die Richtungen auf schon gemachten Messungen gegründet sind, die punktirten projectirte. Ich habe leider Grund zu fürchten, dass der Brelingerberg weder vom Wohlenberg noch vom Lichtenberg sichtbar ist (an allen 3 Orten bin ich selbst nicht gewesen). Überhaupt wird die Gewinnung grosser Dreiecke in der Lüneburger Heide grosse Schwierigkeiten haben.

% Printed 397.
\par\clearpage
\noindent\null\hfill{\small\textbf{ZUR HANNOVERSCHEN TRIANGULATION.}}\hfill\makebox[0pt][r]{\fontsize{13}{16}\selectfont 397}\par\medskip
\par\noindent\hspace*{\parindent}Brelingerberg, Deister, Lichtenberg und Inselsberg sind durch Heliotroplicht sichtbar gemacht, auf der Brockenstation auch Hohehagen und Hils, da die dort gebauten Signale nur selten (letzteres nur wenige Minuten), ich will nicht sagen zu beobachten, sondern nur zu \textls[120]{sehen} gewesen sind. Mit dem Heliotrop fällt alle Schwierigkeit weg.
\par\noindent\hspace*{\parindent}Ich lasse jetzt noch zwei andere machen (nach der neuen Einrichtung), wovon der eine bald vollendet sein wird. Dass man meine telegraphischen mit dem Sextanten-Heliotrop auf dem Brocken gegebenen Zeichen auf dem Hohehagen (Distanz \(70000\) Meter \(=9\frac{1}{3}\) geogr. Meilen) \textls[120]{mit blossen Augen} gesehen, habe ich, wie ich glaube, Ihnen bereits in meinem letzten Briefe gemeldet. Den bisherigen Heliotrop kann mein ältester Sohn Joseph schon recht gut einrichten und lenken; mit dem neuen wird es eher noch etwas leichter gehen.
\par\noindent\hspace*{\parindent}Die Richtung vom Hils auf Hannover ist zwar auch aufs schärfste gemessen; jedoch wird Hannover vermuthlich kein Hauptdreieckspunkt werden, da man von da nach N.O. nur eine sehr begrenzte Aussicht hat, und namentlich den Wohlenberg dort nicht sehen kann. Auch müssten auf dem dortigen Thurme erst grosse Abänderungen gemacht werden, wenn ein Theodolith dort aufgestellt werden sollte. Der Deister wird nach allen Richtungen noch eine ausgedehnte Aussicht beherrschen, und vielleicht wird selbst der Falkenberg da noch gesehen werden können.
\begin{center}
\rule{.74\linewidth}{.4pt}\\[-.5ex]
{\small Astronomische Nachrichten, Band I, Nr.~24, December 1822, S.~441---444.}\\[-.5ex]
\rule{.74\linewidth}{.4pt}\par\medskip
Auszug aus einem Schreiben des Herrn Hofrath \textsc{Gauss} an den Herausgeber.\\
Göttingen 1822, Nov.~10.
\end{center}
\par\noindent\hspace*{\parindent}Da Sie im \(7^{\text{ten}}\) Stück der Astronomischen Nachrichten eine Anzeige über den Stand meiner Triangulirung am Schluss des Jahrs 1821 gegeben haben, so verfehle ich nicht, Ihnen einen kurzen Bericht über den gegenwärtigen Stand der Operationen zu schicken.
\par\noindent\hspace*{\parindent}Im vorigen Jahre waren die fünf Stationen: Göttingen, Meridianzeichen, Hohehagen, Hils, Brocken absolvirt, und vier Punkte für die weitere Fort-

% Printed 398.
\par\clearpage
\noindent\makebox[0pt][l]{\fontsize{13}{16}\selectfont 398}\hfill{\small\textbf{ZUR HANNOVERSCHEN TRIANGULATION.}}\hfill\null\par\medskip
\noindent setzung der Operationen ausgezeichnet, nemlich Lichtenberg, Deister, Wohlenberg und Brelingerberg. Ich fing die Arbeiten des laufenden Jahrs mit einer Recognoscirungsreise in der Lüneburger Heide an, welche ich um so mehr für nothwendig hielt, da ich die grossen Schwierigkeiten, in diesem flachen Lande, welches ohne alle erhebliche Anhöhen und überall schachbrettartig mit Waldung bedeckt ist, ein Dreiecksnetz zu bilden, bereits aus den Berichten des Obersten \textsc{Epailly} kannte, welcher in den Jahren 1804 und 1805 diese Schwierigkeiten unübersteiglich gefunden, und daher die Verbindung zwischen Hamburg und dem südlichen Theile von Hannover vermittelst einer Reihe von Dreiecken längs der Weser bis zu ihrer Mündung und hernach wieder die Elbe herauf effectuirt hatte.
\par\noindent\hspace*{\parindent}Ich fand den Brelingerberg, welcher 1821 vom Hils aus geschnitten war, \textls[120]{unbrauchbar}, da er sich mit Lichtenberg und dem Wohlenberg nicht verbinden liess, aber auch ebenso wie den Wohlenberg \textls[120]{überflüssig}, da sowohl der Platz bei Garssen, als der Falkenberg sich unmittelbar mit dem Lichtenberg verbinden liessen. Ich schweige von den grossen Schwierigkeiten, mit welchen ich zu kämpfen gehabt habe, um die Dreiecke von Garssen und Falkenberg weiter fortzuführen: Diese Schwierigkeiten sind jetzt überwunden, und das Netz bietet durch seinen Gliederbau vielfache zu meiner grössten Zufriedenheit ausgefallene Controllen dar. Ich bemerke nur, dass alle meine Dreieckspunkte zu ebener Erde liegen; ein etwa \(3\frac{1}{2}\)---4 Fuss hoch aufgemauertes steinernes Postament dient zur Aufstellung des Heliotrops und des Theodolithen. Mehrere Linien, namentlich die von Falkenberg nach Wilsede, von Hauselberg nach Breithorn, von Breithorn nach Scharnhorst, und von Scharnhorst [nach Garssen] erforderten beträchtliche Durchhaue durch Waldungen, und die genaue Vorausbestimmung der Richtung dieser Durchhaue künstliche Vorbereitungen.
\par\noindent\hspace*{\parindent}Ich habe im Laufe des Sommers die Stationen Lichtenberg, Deister, Garssen, Falkenberg, Hauselberg, Breithorn, Wulfsode, Wilsede und Scharnhorst vollständig abgemacht, auch auf Timpenberg die betreffenden Winkel vorläufig gemessen. Dadurch ist also Hamburg schon vorläufig angeschlossen, auch Lüneburg, da Sie den Winkel zwischen Wilsede und Lüneburg auf dem Michaelisthurm in Hamburg vorläufig gemessen haben. Hier einige vorläufige Resultate, wobei sich das Absolute vorläufig auf die von \textsc{Zach}sche Basis bei

% Printed 399.
\par\clearpage
\noindent\null\hfill{\small\textbf{ZUR HANNOVERSCHEN TRIANGULATION.}}\hfill\makebox[0pt][r]{\fontsize{13}{16}\selectfont 399}\par\medskip
\noindent Gotha gründet, an die ich mich vermittelst der Seite vom Inselsberg zum Brocken angeschlossen habe.
\begin{center}
\setlength{\tabcolsep}{3pt}
\begin{tabular}{@{}lrrr@{\qquad}rrrl@{}}
 & \multicolumn{3}{c}{Breite} & \multicolumn{4}{c}{\shortstack{Länge von\\Göttingen}}\\[-1.5ex]
 & \multicolumn{3}{c}{\(\overbrace{\hphantom{53\quad33\quad31{,}4}}\)} & \multicolumn{4}{c}{\(\overbrace{\hphantom{0\quad27\quad29{,}5\quad\text{östl.}}}\)}\\[-1.2ex]
Hamburg, Michaelisthurm & \(53^\circ\) & \(33'\) & \(1^{\prime\prime}{,}8\) & \(0^\circ\) & \(2'\) & \(3^{\prime\prime}{,}0\) & östl.\\
Lüneburg, Michaelisthurm & 53 & 15 & \(5{,}5\) & 0 & 27 & \(29{,}5\) & \multicolumn{1}{c}{»}\\
Celle, südl. Schlossthurm & 52 & 37 & \(31{,}4\) & 0 & 8 & \(4{,}9\) & \multicolumn{1}{c}{»}\\
Göttinger Sternwarte, Platz & & & & & & &\\
\quad des \textsc{Reichenbach}schen & & & & & & &\\
\quad Meridiankreises & 51 & 31 & \(48{,}7\) & & 0 & &
\end{tabular}
\end{center}
\par\noindent\hspace*{\parindent}Die Orientirung meines Dreieckssystems ist von meinem Meridianzeichen entlehnt; auf Hamburg übertragen weicht sie von den Azimuthen, welche Sie mir mitgetheilt haben, nur \(1^{\prime\prime}{,}4\) ab. Um das Absolute der Linien schärfer zu bestimmen, erwarte ich nur die Mittheilung der Länge Ihrer Basis und die Dreiecke, welche sie mit Ihren Hauptpunkten verbindet. Meine eine Dreiecksseite, Breithorn-Scharnhorst, würde sich, wie es scheint, ohne unübersteigliche Schwierigkeiten unmittelbar messen lassen.
\par\noindent\hspace*{\parindent}Um eine recht zweckmässige Verbindung meiner Dreiecke mit den Ihrigen zu erhalten, hatte ich gewünscht und gehofft, Timpenberg mit Lüneburg unmittelbar verbinden zu können. Ein Durchhau wurde versucht, allein, nachdem er eine bedeutende Strecke hindurch fortgeführt war, fand sich schon das zwischenliegende Terrain nicht deprimirt genug, und musste daher diese unmittelbare Verbindung aufgegeben werden. Es ist jedoch von Wilsede aus noch ein Punkt niedergelegt, der sich unmittelbar mit Hamburg, Lüneburg und Lauenburg und höchst wahrscheinlich vermittelst eines Durchhaues mit Timpenberg verbinden lassen wird. Das Weitere muss den Arbeiten des künftigen Jahrs vorbehalten bleiben. Von den grossen Schwierigkeiten, in einem solchen waldigen, flachen Terrain zu operiren, hat Niemand einen Begriff, der nicht unter ähnlichen Umständen gearbeitet hat.
\par\noindent\hspace*{\parindent}Die beifolgende Karte, welche in dem Maassstabe von \(\frac{1}{2000000}\) gezeichnet ist, wird Ihnen von dem Geschafften eine anschaulichere Vorstellung geben. Erst nachdem die übrigen Arbeiten vollendet waren, fand sich, dass der Punkt Scharnhorst, vermittelst zweier nicht sehr schwieriger Durchhaue, unmittelbar

% Printed 400.
\par\clearpage
\noindent\makebox[0pt][l]{\fontsize{13}{16}\selectfont 400}\hfill{\small\textbf{ZUR HANNOVERSCHEN TRIANGULATION.}}\hfill\null\par\medskip
\noindent mit Lichtenberg und Deister sich verbinden lassen würde. Wäre es möglich gewesen, diesen Platz früher auszumitteln, und seine Brauchbarkeit und Lage festzusetzen, so hätte Garssen ganz wegfallen können. Vielleicht werde ich im künftigen Jahre die Messung der Winkel des Dreiecks Scharnhorst-Deister-Lichtenberg noch nachholen.
\par\noindent\hspace*{\parindent}Ich habe in diesem Jahre ausser dem im Jahre 1821 gebrauchten Heliotrop noch zwei andere von der neuen Einrichtung in Thätigkeit gehabt, und daneben noch einen andern Heliotrop-Apparat, welchen ich immer bei mir führte, um meinen Gehülfen telegraphische Ordres zu geben. Für Sie ist die Bemerkung überflüssig, dass die von Hrn.~\textsc{Schubach} im Astron. Jahrbuch für 1825 gegebene Nachricht über die Einrichtung der Heliotrope ganz auf einem Irrthum beruht und mit \textls[120]{meinen} Heliotropen gar nichts gemein hat. Herr \textsc{Rumpf} hat bereits sieben Heliotrope verfertigt, wovon zwei für die preussische und zwei für die hessische Triangulirung bestimmt sind. Von beiden Einrichtungen stehen Ihnen auf Verlangen Zeichnungen zu Dienste.
\begin{center}\rule{.12\linewidth}{.4pt}\end{center}
\begin{center}\textbf{BEMERKUNGEN.}\end{center}
\begingroup\small
\par\noindent\hspace*{\parindent}Die dem ersten Artikel in Band I, Nr.~7, der Astronomischen Nachrichten, beigegebene Dreieckskarte enthält in starken Linien im Süden die Dreiecke von der Seite Hohehagen-Inselsberg bis zur Seite Hils-Lichtenberg (vergl. die Figur auf S.~299), die Schnitte vom Hils zum Deister und Brelingerberg und im Norden die beiden Dreiecke Lüneburg-Hamburg-Hohenhorn und Lüneburg-Hohenhorn-Lauenburg. Durch punktirte Linien sind gebildet das Dreieck Lichtenberg-Deister-Wohlenberg, sowie die Richtungen zwischen den Punkten Deister, Brelingerberg, Falkenberg und Wilsede und die Richtungen Wilsede-Lüneburg und Wilsede-Hamburg. Wohl aus Versehen fehlt auf dieser Karte die Linie Hils-Nördl. Meridianzeichen.
\par\noindent\hspace*{\parindent}Die dem zweiten Artikel zugefügte Dreieckskarte gibt die Dreiecke der eigentlichen hannoverschen Gradmessung, also die in der Figur auf S.~299 dargestellten Dreiecke vom Inselsberg bis Hamburg bis zu den Seiten Falkenberg-Wilsede und Wilsede-Hamburg. Es fehlt aber auf ihr noch der Punkt Nindorf und seine Verbindungen; dagegen sind auf ihr die Anschlussdreiecke zwischen den Punkten Lüneburg, Wilsede, Hamburg und Hohenhorn und die beiden dänischen Punkte Syk und Lauenburg enthalten.
\par\noindent\hspace*{\parindent}Ausser diesen beiden Mittheilungen über die hannoversche Triangulirung sind von \textsc{Gauss} in Beziehung auf dieselbe noch 2 (später folgende) Aufsätze aus dem Jahre 1821 über den »Heliotrop« und ein (gleichfalls noch folgender) Artikel aus dem Jahre 1823: »Beobachtete und berechnete Triangulirung im Hannoverschen, Braunschweigschen und Lüneburgschen« veröffentlicht worden.
\par\noindent\hfill\textsc{Krüger}.\par
\endgroup
\begin{center}\rule{.16\linewidth}{.4pt}\end{center}
% End of printed 400. Printed 401 opens III. NACHLASS ZUR HANNOVERSCHEN TRIANGULATION.
